\subsection{Factorización a través del operador constitutivo}

Conservemos la coordenatización angular levantada de la
sección~\ref{iii:sec:hojas}. Para hacer explícita la conversión angular,
escribamos
\begin{equation}
 A_{\deg}=\frac{180x}{\pi},\qquad
 C^*_{\deg}=\frac{180y}{\pi},\qquad
 q_\pm=e^{-x\mp y},\qquad s_\pm=q_\pm^{30}.
 \label{iii:eq:barbero-angular-lift}
\end{equation}
La orientación considerada en \eqref{iii:eq:rchannels} corresponde a
$0<y<x$; el intercambio de hojas se expresa por $y\mapsto-y$ en el
dominio más amplio $x>|y|$. Las exponenciales reales actúan sobre estos
levantamientos, que conservan los valores de $x\pm y$ durante la
composición.

El funcional orientado de Barbero--Immirzi es
\begin{equation}
 \gamma=\frac{\pi A_{\deg}C^*_{\deg}}{180}
   +12\bigl[S_{90}(q_-)-S_{90}(q_+)\bigr]
   -\bigl[S_{120}(q_-)-S_{120}(q_+)\bigr].
 \label{iii:eq:barbero-functional}
\end{equation}
La primera componente es la evaluación angular; la segunda es la
diferencia del lector de la cadena fundamental en los canales conjugados.
Los argumentos se reciben del transporte angular previamente construido.
La fórmula determina un escalar adimensional sobre esa estructura de
partida.

Sea $\psi:(3/4,1)\to(0,1)$ la inversa de $f$ construida mediante
\eqref{iii:eq:cubic}. Definamos
\begin{equation}
 \mathcal H(s)=\frac{12s^3}{1-s^9}
             -\frac{s^4}{1-s^{12}}
             -\frac{(\log s)^2}{20\pi},\qquad 0<s<1.
 \label{iii:eq:barbero-h}
\end{equation}
Esta función es real analítica en su dominio. Los exponentes $3,4,9,12$
resultan de expresar las alturas $90,120$ y los retornos $270,360$ en
unidades de $30$; el coeficiente de la parte logarítmica se fija al
transportar la contribución angular.

\begin{theorem}[Expresión constitutiva orientada de Barbero--Immirzi]
\label{iii:thm:barbero-factorization}
En la coordenatización \eqref{iii:eq:barbero-angular-lift},
\begin{equation}
 \gamma=\mathcal H(s_-)-\mathcal H(s_+)
       =\mathcal H\bigl(\psi(r_-)\bigr)
        -\mathcal H\bigl(\psi(r_+)\bigr).
 \label{iii:eq:barbero-factorization}
\end{equation}
Sobre la representación bidimensional de canales,
\begin{equation}
 \gamma=-\operatorname{Tr}
    \left[\mathcal R\,\mathcal H\bigl(\psi(\mathcal V)\bigr)\right],
 \label{iii:eq:barbero-trace}
\end{equation}
donde $\operatorname{Tr}$ es la traza ordinaria, con
$\operatorname{Tr}I=2$.
\end{theorem}
\begin{proof}
La sustitución $s=q^{30}$ en \eqref{iii:eq:barbero-return-series} da
\[
 S_{90}(q)=\frac{s^3}{1-s^9},\qquad
 S_{120}(q)=\frac{s^4}{1-s^{12}}.
\]
Además, \eqref{iii:eq:barbero-angular-lift} implica
\[
 \log s_\pm=-30(x\pm y)
 =-\frac\pi6(A_{\deg}\pm C^*_{\deg}).
\]
Por diferencia de cuadrados,
\begin{align}
 \frac{(\log s_+)^2-(\log s_-)^2}{20\pi}
 &=\frac{900\bigl[(x+y)^2-(x-y)^2\bigr]}{20\pi}\\
 &=\frac{180xy}{\pi}
 =\frac{\pi A_{\deg}C^*_{\deg}}{180}.
 \label{iii:eq:barbero-log-term}
\end{align}
La suma de estas identidades prueba la primera igualdad de
\eqref{iii:eq:barbero-factorization}. La segunda utiliza la inversión
$\psi(f(s_\pm))=s_\pm$ demostrada en la
sección~\ref{iii:sec:constitutiva}.

Para la expresión operatoria, los proyectores de hoja tienen rango uno y
$\mathcal RP_\pm=\pm P_\pm$. El cálculo espectral proporciona
\[
 \mathcal H\bigl(\psi(\mathcal V)\bigr)
 =\mathcal H(s_+)P_++\mathcal H(s_-)P_-.
\]
Multiplicar por $\mathcal R$ y tomar la traza ordinaria da
$\mathcal H(s_+)-\mathcal H(s_-)$. Su opuesto es $\gamma$.
\end{proof}

La marca de orientación $\mathcal R$ forma parte del funcional. Si se
mantiene fija esa marca y se intercambian los canales de $\mathcal V$,
cambia el signo de $\gamma$. La evaluación del espectro sin sus etiquetas conserva
el par no ordenado, mientras \eqref{iii:eq:barbero-trace} conserva la
contribución orientada. Si se emplea la traza normalizada
$\tau_2=\operatorname{Tr}/2$, la misma fórmula se escribe
$\gamma=-2\tau_2[\mathcal R\mathcal H(\psi(\mathcal V))]$.

En cambio, un cambio simultáneo de representación
$(\mathcal R,\mathcal V)\mapsto
(U\mathcal RU^{-1},U\mathcal VU^{-1})$ conserva el funcional:
el cálculo espectral se transporta por conjugación y la traza es
invariante. Intercambiar la respuesta respecto de una orientación fija
y transportar conjuntamente respuesta y orientación son, por tanto,
operaciones diferentes.

Esta normalización es compatible con las amplificaciones de
\eqref{iii:eq:logvac}. Para
$\mathcal V_m=\mathcal V\otimes I_{9^m}$ y
$\mathcal R_m=\mathcal R\otimes I_{9^m}$,
\begin{equation}
 \gamma=-\frac1{9^m}\operatorname{Tr}
 \left[\mathcal R_m\mathcal H\bigl(\psi(\mathcal V_m)\bigr)\right].
 \label{iii:eq:barbero-amplification}
\end{equation}
En efecto, el cálculo funcional conserva el producto tensorial y la traza
de $I_{9^m}$ es $9^m$. La multiplicidad de la representación se separa
de la evaluación orientada.

\subsection{Reconstrucción a partir de las coordenadas del vacío}

Las identidades \eqref{iii:eq:four-identities} permiten expresar el mismo
funcional mediante la impedancia y la velocidad normalizadas:
\begin{align}
 \gamma={}&
 \mathcal H\!\left(\psi\!\left(
     \sqrt{\frac{\widehat Z}{\widehat c}}\right)\right)
 -\mathcal H\!\left(\psi\!\left(
     \frac1{\sqrt{\widehat Z\widehat c}}\right)\right).
 \label{iii:eq:barbero-zc}
\end{align}
El dominio positivo de esta reconstrucción se describe por
\begin{equation}
 1<\widehat Z\widehat c<\frac{16}{9},\qquad
 \frac9{16}<\frac{\widehat Z}{\widehat c}<1.
 \label{iii:eq:barbero-zc-domain}
\end{equation}
En efecto, esas desigualdades equivalen a
$3/4<r_\pm<1$ al sustituir las raíces positivas de
\eqref{iii:eq:four-identities}. Para la orientación $y>0$ se conserva
además $r_-<r_+$, equivalentemente $\widehat Z<1$. Los puntos del
dominio se utilizan con la compatibilidad angular ya establecida por
el TPK; la inversión es una reconstrucción posterior de sus canales.
Las coordenadas son las internas de \eqref{iii:eq:four}: un cambio de
unidades debe deshacerse mediante \eqref{iii:eq:undo-units} antes de
aplicar $\psi$.

El intercambio $y\mapsto-y$ actúa como
$(\widehat Z,\widehat c)\mapsto(\widehat Z^{-1},\widehat c)$ y
produce $\gamma\mapsto-\gamma$. En el caso balanceado $y=0$,
ambos canales coinciden y \eqref{iii:eq:barbero-factorization} da
$\gamma=0$. La velocidad normalizada conserva el producto de los
canales; su lectura conjunta con la impedancia recupera las dos
respuestas y la orientación necesaria para evaluar el funcional.
En ese caso balanceado $\mathcal V=rI$ tiene espectro degenerado:
la marca $\mathcal R$ se conserva como dato de representación, aunque
el escalar $r$ ya no distinga sus proyectores. Para espectro simple,
en cambio,
$P_+=(\mathcal V-r_-I)/(r_+-r_-)$ y $P_-=I-P_+$.
